\subsection{Выбор средств реализации}
\label{sec:development:tools}

Средства реализации выбирались под три особенности задачи. Данные изменяются одновременно несколькими процессами. Большую часть времени система ждёт ответов от внешних сервисов: площадки, \textit{Telegram}, почтового сервера. Система работает круглосуточно на одном сервере и должна переживать перезапуски без потери данных.

\subsubsection{Язык программирования}

Система написана на \textit{Python}~3.11. В языке есть встроенная поддержка асинхронного выполнения, при котором один процесс обслуживает множество одновременных операций ввода-вывода, не создавая поток на каждую из них. Пока система ждёт ответа площадки на один запрос, параллельно обрабатываются команды пользователей и отправляются уведомления. Для \textit{Python} существуют зрелые асинхронные библиотеки для всех внешних сервисов системы: \textit{Telegram}, \textit{HTTP}, \textit{PostgreSQL}.

\subsubsection{Выбор СУБД}

В качестве СУБД выбрана \textit{PostgreSQL}~16. Наряду с ней рассматривались \textit{MySQL}, встраиваемая \textit{SQLite} и документоориентированная \textit{MongoDB}. Сравнение по свойствам, значимым для задачи, приведено в таблице~\ref{table:development:dbms} [7].

\begin{table}[ht]
\caption{Сравнение СУБД по требованиям системы}
\label{table:development:dbms}
\centering
	\begin{tabular}{|>{\raggedright}m{0.32\textwidth}
	                |>{\centering}m{0.13\textwidth}
	                |>{\centering}m{0.13\textwidth}
	                |>{\centering}m{0.13\textwidth}
	                |>{\centering\arraybackslash}m{0.13\textwidth}|}
	\hline
	\centering Свойство & \textit{Post\-greSQL} & \textit{MySQL} & \textit{SQLite} & \textit{Mongo\-DB} \\
	\hline Работа как отдельный сетевой сервер & $+$ & $+$ & $-$ & $+$ \\
	\hline Одновременная запись несколькими процессами & $+$ & $+$ & $-$ & $+$ \\
	\hline Внешние ключи & $+$ & $+$ & $+$ & $-$ \\
	\hline Откат изменений схемы в транзакции & $+$ & $-$ & $+$ & не требуется \\
	\hline Хранение \textit{JSON}-документов & $+$ & $+$ & частично & $+$ \\
	\hline Рекомендательные блокировки & $+$ & $+$ & $-$ & $-$ \\
	\hline
	\end{tabular}
\end{table}

\textit{SQLite} хранит базу в одном файле и в каждый момент допускает только одну пишущую транзакцию. Первая версия бота работала на ней в одном процессе. Когда система была разделена на несколько сервисов, работающих с общими данными, от встраиваемой базы пришлось отказаться в пользу сетевой СУБД. \textit{MongoDB} не поддерживает внешние ключи, а связи между пользователем, товаром, совпадением и уведомлением составляют основу модели данных, и их целостность пришлось бы контролировать вручную. \textit{MySQL} подошла бы по большинству свойств, однако изменение схемы в ней не откатывается вместе с транзакцией: миграция, прерванная на середине, оставляет базу в промежуточном состоянии.

\textit{PostgreSQL} закрывает все требования. Внешние ключи и транзакции обеспечивают целостность цепочки от пользователя до уведомления. Тип \textit{JSON} хранит данные переменной структуры: фильтры площадки, у которых набор характеристик свой для каждой категории, списки слов, адреса фотографий, исходный ответ площадки. Выделять под них отдельные таблицы не нужно. Рекомендательная блокировка (\textit{advisory lock}) защищает от одновременного применения миграций несколькими экземплярами сервиса при запуске.

\subsubsection{Доступ к базе данных и миграции}

Запросы к базе данных записываются на SQL и выполняются через асинхронный драйвер \textit{asyncpg}, который работает с сетевым протоколом \textit{PostgreSQL} напрямую. Структура базы данных описывается версионированными миграциями под управлением \textit{Alembic}. Каждая миграция описывается SQL-скриптом перехода между двумя версиями схемы и содержит обратный переход. Миграции применяются автоматически при запуске серверного сервиса, поэтому схема рабочей базы всегда соответствует версии кода [8, 9].

\subsubsection{Серверный интерфейс}

Программный интерфейс, через который бот получает и изменяет данные, построен на веб-фреймворке \textit{FastAPI} и запускается сервером \textit{uvicorn}. \textit{FastAPI} асинхронный и проверяет входные данные по объявленным схемам \textit{pydantic}: запрос с отрицательной ценой или слишком длинным названием отклоняется до обращения к базе данных. Описание интерфейса формируется автоматически по коду. \textit{Flask} требует для тех же возможностей набора расширений, а \textit{Django} рассчитан на веб-приложения с шаблонами страниц и административной панелью, которые системе не нужны [10].

\subsubsection{Работа с \textit{Telegram}}

Бот построен на библиотеке \textit{aiogram}~3. Библиотека асинхронная и поддерживает конечные автоматы для многошаговых диалогов: добавление товара проходит до восьми шагов, и на каждом шаге сохраняется состояние диалога пользователя. Обработчики разделяются по маршрутизаторам, а общая логика, например определение пользователя по каждому входящему сообщению, выносится в промежуточный слой [11].

\subsubsection{Периодические задачи}

Опрос площадки, доставка уведомлений, рассылка почтовых копий, очистка устаревших данных и обновление справочника фильтров запускаются планировщиком \textit{APScheduler} внутри серверного сервиса. Отдельный брокер сообщений, как у \textit{Celery}, для пяти задач на одном сервере не требуется, а системный \textit{cron} не даёт запретить повторный запуск задачи, пока предыдущий ещё не завершился [12].

\subsubsection{Вспомогательные библиотеки}

Запросы к площадке выполняет асинхронный клиент \textit{aiohttp}, бот обращается к серверному интерфейсу через \textit{httpx}. Нечёткое сравнение строк при проверке названия товара выполняет \textit{rapidfuzz}. Отчёты в формате \textit{Excel} формирует \textit{openpyxl}. Журнал работы с ротацией файлов ведёт \textit{loguru}. Тесты написаны на \textit{pytest}.

\subsubsection{Развёртывание}

Все компоненты запускаются в контейнерах \textit{Docker} и описаны одним файлом \textit{Docker Compose}: СУБД, серверный сервис, бот, резервное копирование, служба мониторинга контейнеров и почтовый сервер для проверки писем. Такой способ воспроизводит окружение на любом сервере одной командой, автоматически перезапускает упавшие контейнеры и хранит данные базы в отдельном томе, который не пропадает при пересборке [13].
